\chapter{Reliability over Astronomical Intervals}
\label{ch:29}

The most ambitious objective in the program is the extension of human life beyond centuries and millennia toward durations measured in millions or billions of years. The objective should be stated with exactness about what is and is not presently established. A million-year lifespan is not a demonstrated medical possibility, and a billion-year lifespan raises further astronomical and probabilistic difficulties, among them the fact that the interval is comparable to the remaining period over which the Earth is expected to remain habitable. The approach begins with biological aging, cellular repair, the prevention of cancer, regenerative medicine, the accumulation of molecular damage, and the limits of organismal maintenance. It then considers more speculative possibilities, including comprehensive tissue replacement, synthetic biological systems, and methods of preserving and restoring neural structure, with the explicit understanding that whether any of these maintains personal continuity is a question on which scientific and philosophical analysis remain unsettled, and that the answer may differ among approaches.

The essential observation is that extreme longevity is not only a problem of aging. Even if senescence were eliminated entirely, a person would remain exposed to accident, violence, infectious and environmental hazard, institutional failure, and the cumulative effect of small risks that are negligible over a decade and decisive over a million years. The mathematics of this exposure imposes quantitative requirements that are independent of any particular technology. Suppose the annual probability of fatal failure is a constant $p$ and that years are independent. Survival to time $T$ then has probability $(1-p)^T$, which for small $p$ is well approximated by $e^{-pT}$. For survival over a million years with probability near one-half, the annual risk would have to be below a figure on the order of one in a million, and for a billion years below a figure on the order of one in a billion. For comparison, the annual mortality of healthy young adults in high-income countries is on the order of one in a thousand, so that the required reduction in risk is three orders of magnitude for the first case and six for the second, and it must be sustained and not merely achieved. Reductions of this size cannot be obtained by treating disease alone, and they require the management of every class of hazard, together with the capacity to repair damage faster than it accumulates and to maintain redundancy against failures that cannot be repaired.

The mathematics also exposes a limit on redundancy that is easily missed. Maintaining multiple independent backups can drive the probability of simultaneous failure very low, but the argument fails when failures share a common cause, such as an environmental catastrophe, a flaw in a shared design, or a failure of the institutions that maintain the system, and the probability of such common-cause failure imposes a floor that no degree of replication lowers. For this reason the question of longevity is inseparable from the questions of institutional durability and environmental stability that other parts of the program address. The final reflection concerns what becomes of education, law, inheritance, population, memory, institutions, and identity when individual lives might extend across geological epochs. Legal and social structures presuppose generational turnover, and a population in which individuals persist indefinitely would require new provision for the renewal of authority, the redistribution of resources, the revision of commitments made long ago, and the preservation of the capacity to change one's mind and one's way of life. These questions are not consequences to be deferred until the biology is solved; they bear on whether the objective is desirable in the form in which it is stated, and they belong among the admissibility constraints of the program.

The sources of the hazard deserve to be listed, since the requirement of the preceding paragraph is a requirement on their sum. Some are natural and astronomical. Impacts of asteroids and comets large enough to cause global catastrophe, of the order of a kilometer in diameter, are estimated to occur on average once in a few hundred thousand years to a million, and the impact that ended the Cretaceous period, some sixty-six million years ago, was larger. Volcanic eruptions of the supervolcanic class, which can alter the climate of the whole planet for years, recur at intervals of the order of hundreds of thousands of years, and the eruption of Toba some seventy-four thousand years ago is the best known. Nearby supernovae and gamma-ray bursts are rare but not negligible over geological time. On the longest scale the luminosity of the Sun rises by about one per cent in a hundred million years, and the Earth is expected to lose its oceans in a time that is estimated to be of the order of a billion years, which means that the aim of a billion-year life cannot be met on the Earth alone, whatever is done about every other hazard. Other hazards are of human origin: war, and in particular nuclear war; engineered or natural pandemics; the failure of the complex systems on which life depends; and the actions of a person or group intent on harm. The first class has an irreducible rate on the surface of a single planet, to be reduced only by dispersal; the second has a rate that depends on institutions and on the choices that societies make, and it is the one on which the program can act soonest.

The consequence for the design of a long-lived civilization is that its hazard must be reduced by structure and not merely by vigilance. The most powerful structural measure is dispersal, the maintenance of populations in locations whose catastrophes are not correlated, so that the loss of any one does not end the whole. If $k$ sites fail independently with annual probability $p_s$ each, the probability that all fail in a given year is $p_s^k$, and the survival of a civilization over the age of the Earth becomes plausible with a modest $k$ provided that the sites are truly independent, which, as the beta-factor analysis below shows, they never are entirely. The common causes are the shared artifacts of human origin, the conflicts that arise among the sites, and the dependence of each on the same institutions, so that the floor on the failure probability is set by the weakest of the cross-site links. This is the point at which the orbital programs of the fifth part and the longevity program of the sixth meet: settlements in space, if they could be maintained, would be the sites whose failures are least correlated with those of the Earth, and the case for the orbital programs, apart from their direct benefits, is that they are the means by which the hazard of the whole can be brought toward the required values. The case is conditional and is the strongest form of the longtermist argument, which Chapter 31 examines; it can be accepted as a reason to keep the options open and to prepare, and it cannot be accepted as a reason to override present obligations.

A second structural measure is the detection and response to hazards before they arrive. The survey of near-Earth objects, which has already cataloged most objects of a kilometer in size and found none on a collision course in the coming century, and the demonstration of a deflection technique by an impact on a small asteroid in 2022, are instances in which a natural hazard has been moved from the class of the unavoidable to that of the manageable. The same logic applies to pandemics, with surveillance and rapid response, and to the vulnerability of the systems of energy, communication, and food. Each reduction is modest, a factor of two or ten in a hazard that must be reduced by factors of millions, and the lesson of the calculation that follows is that a long life is the product of many such reductions, each of which must be maintained without lapse. The requirement of independence between verifier and generator, and the readiness classification, are in this sense not peripheral to the longevity program but constitutive of it: an institution that cannot maintain its own reliability over a century cannot be the guarantor of a person's survival over a million years.

\section*{Formalization}

Let $h:[0,\infty)\to[0,\infty)$ be the hazard rate of fatal failure at time $t$ (in years), with cumulative hazard $H(T)=\int_0^T h(t)\,dt$. The survival function is $S(T)=\exp(-H(T))$. For a constant discrete annual risk $p$, $S(T)=(1-p)^T=\exp\big(T\ln(1-p)\big)$.

\begin{proposition}
For constant annual risk $p\in(0,1)$ and a survival target $q\in(0,1)$ over $T$ years, the required annual risk is $p^\ast=1-q^{1/T}$, and $p^\ast=-\ln q/T+O(T^{-2})$. In particular, for $q=\tfrac12$, $p^\ast\approx(\ln2)/T\approx0.693/T$.
\end{proposition}

\begin{proof}
Solving $(1-p)^T=q$ gives $p=1-q^{1/T}$. Writing $q^{1/T}=\exp\big((\ln q)/T\big)=1+(\ln q)/T+O(T^{-2})$ gives the expansion, and $\ln\tfrac12=-\ln2$.
\end{proof}

The following table gives $S(T)=(1-p)^T\approx e^{-pT}$ for the stated values.

\begin{center}
\begin{tabular}{@{}lcc@{}}
\toprule
Annual fatal risk $p$ & $T=10^{6}$ years & $T=10^{9}$ years \\
\midrule
$10^{-3}$  & $\approx e^{-1000}$ (negligible) & $\approx e^{-10^{6}}$ (negligible) \\
$10^{-6}$  & $\approx 0.368$ & $\approx e^{-1000}$ (negligible) \\
$10^{-9}$  & $\approx 0.999$ & $\approx 0.368$ \\
$10^{-12}$ & $\approx 0.999999$ & $\approx 0.999$ \\
\bottomrule
\end{tabular}
\end{center}

Here $e^{-1}\approx0.368$ occurs whenever $pT=1$, which marks the characteristic scale: a risk $p$ supports survival over about $1/p$ years with probability $1/e$.

For redundancy consider a single repair interval and $k$ replicas, under the following explicit assumptions: a common-cause event occurs with probability $q_c$ and fails every replica simultaneously; independently of it, each replica suffers an independent failure with probability $q_i$, independently across replicas; and failure of all replicas is the event of interest. Then
\[
p_{\mathrm{eff}}=q_c+(1-q_c)\,q_i^{\,k}.
\]
In the beta-factor parameterization, in which $p_f$ denotes the per-replica failure probability (taken small, so that it is approximately $q_c+q_i$) and $\beta\in[0,1]$ the common-cause fraction, $q_c=\beta p_f$ and $q_i=(1-\beta)p_f$, so that
\[
p_{\mathrm{eff}}=\beta p_f+(1-\beta p_f)\big((1-\beta)p_f\big)^{k}\;\ge\;\beta\,p_f.
\]
Increasing $k$ drives the second term to zero but leaves the floor $\beta p_f$ per repair interval, which accumulates over the intervals contained in a year. The achievable annual risk is therefore bounded below by a quantity proportional to $\beta p_f$, and the requirement $p^\ast$ of the proposition can be met only if that floor does not exceed it, a condition on the correlation of failures that replication cannot substitute for. Over a population of $N$ individuals the expected number of survivors is $N\,S(T)$, and if the cumulative hazard $H$ varies across individuals the expectation of survival is $\E[e^{-H}]\ge e^{-\E[H]}$ by Jensen's inequality.
